% © 2025 Ing. David Jaroš — CC BY-NC-ND 4.0
%
% This work is licensed under a Creative Commons Attribution-NonCommercial-NoDerivatives
% 4.0 International License (CC BY-NC-ND 4.0).
%
% License History: Earlier drafts (up to v0.3) were released under CC BY 4.0.
% From v0.4 onward, all material is released under CC BY-NC-ND 4.0 to protect
% the integrity of the theoretical work during ongoing academic development.

\documentclass[12pt]{article}
\usepackage{amsmath,amssymb,amsthm}
\usepackage{geometry}
\usepackage{hyperref}
\usepackage{tcolorbox}
\usepackage{xcolor}
\geometry{margin=1in}

% Theorem-like environments
\newtheorem{conjecture}{Conjecture}
\newtheorem{proposition}{Proposition}
\newtheorem{definition}{Definition}
\newtheorem{remark}{Remark}

\title{ST-5: Methodology Summary —\\
Closed Timelike Curves and Time Travel in\\
Unified Biquaternion Theory\\[0.5em]
\large\textsc{Speculative Extension}}
\author{David Jaroš}
\date{2026-03-01}

\begin{document}
\maketitle

\begin{tcolorbox}[colback=yellow!10, title=Status]
\textbf{SPECULATIVE EXTENSION} — This content is not part of
the canonical UBT framework. It is an exploratory investigation
and has not been validated against the canonical field equations.
\end{tcolorbox}

\vspace{1em}
\begin{tcolorbox}[colback=red!5, title=Framing Requirement]
All claims in this document are \textbf{SPECULATIVE}.
No claim asserts that time travel is possible.
All statements are conditional: they describe what field configurations
\emph{would be necessary if} CTCs exist within UBT.
\end{tcolorbox}

\vspace{0.5em}
\noindent\textit{Output of Sub-task ST-5 (methodology summary) from the
CTC/Causal-Structure investigation track.  This document synthesises the
results of ST-1 through ST-4:
\texttt{complex\_time\_causal\_structure.tex},
\texttt{rotation\_complex\_time.tex},
\texttt{ctc\_conditions.tex},
\texttt{chronology\_protection\_ubt.tex}.}

%----------------------------------------------------------------------
\section*{Abstract}
%----------------------------------------------------------------------
We synthesise the findings of sub-tasks ST-1 through ST-4 into a concise
methodology for investigating closed timelike curves (CTCs) in the
Unified Biquaternion Theory (UBT).  We identify the necessary conditions
for CTC existence, the minimal modification to the canonical field
equations that could admit CTCs, and describe what a ``time-travel
protocol'' would look like in terms of field configurations of the UBT
master field $\Theta$.  All claims are framed as conjectures and are
strictly conditional on listed assumptions.

%----------------------------------------------------------------------
\section{Synthesis of ST-1 through ST-4}
%----------------------------------------------------------------------
\subsection{Causal Structure (ST-1)}
Complex time $\tau = t + i\psi$ admits three interpretations of $\psi$:
\begin{itemize}
  \item \textbf{B (preferred):} $\psi$ is an internal phase coordinate.
    Causal order is determined by real time $t$; no new causal pathologies.
  \item \textbf{A (speculative):} $\psi$ is a second independent time
    direction; closed paths in $(t,\psi)$ can wrap around $\psi$ and
    return to the same point while $t$ advances — a direct CTC mechanism.
  \item \textbf{C (degenerate):} $\psi$ is auxiliary; no causal implications.
\end{itemize}
The canonical UBT framework favours Interpretation B.
Interpretation A is required for the most direct CTC mechanism.

\subsection{Complex-Time Rotations (ST-2)}
The $SO(2)$ rotation $\tau \mapsto e^{i\alpha}\tau$ in the complex time
plane:
\begin{itemize}
  \item $\alpha = \pi/2$ (Wick rotation): analytic continuation, not a
    physical symmetry.
  \item General $\alpha$: possible $U(1)_\tau$ symmetry if the UBT action
    is invariant; generates a conserved temporal charge $Q_\tau$.
  \item $\alpha = \pi$ (full inversion): candidate for $T$-reversal in UBT;
    combined with $C$ and $P$, extends the CPT symmetry.
\end{itemize}
The existence of a $U(1)_\tau$ symmetry would generate a topological
obstruction to CTC formation: $Q_\tau$ must be conserved along any closed
path, which constrains which closed paths are physically accessible.

\subsection{CTC Conditions (ST-3)}
The biquaternionic metric $\mathcal{G}_{\mu\nu}$ contributes an effective
stress-energy $T_{\mu\nu}^{\mathrm{eff}}$ to the real Einstein equation.
The imaginary sector $h_{\mu\nu}$ can provide negative effective energy
density, which:
\begin{itemize}
  \item Is a necessary ingredient for CTC-enabling geometries in GR
    (energy condition violation).
  \item In UBT arises from the $\Theta$-field itself without ad hoc exotic
    matter.
  \item Both the Gödel and Kerr metrics can be embedded as real projections
    of UBT solutions with non-trivial $h_{\mu\nu}$.
\end{itemize}

\subsection{Chronology Protection (ST-4)}
Hawking's chronology protection conjecture rests on QFT assumptions that are
modified in UBT:
\begin{itemize}
  \item Semiclassical gravity is replaced by coupled biquaternionic dynamics.
  \item The $\psi$ sector provides a UV regulator for high-frequency modes.
  \item The Cauchy-horizon divergence may be softened (finite KK sum) or
    enhanced (divergent KK sum), corresponding to chronology protection
    circumvention or enforcement, respectively.
\end{itemize}

%----------------------------------------------------------------------
\section{Necessary and Sufficient Conditions for CTCs in UBT}
%----------------------------------------------------------------------
Based on the synthesis above, we state the following:

\begin{conjecture}[Necessary conditions for CTCs in UBT]
\label{conj:necessary}
\textbf{(Speculative)}  If a UBT solution admits a CTC, then all of the
following hold:
\begin{enumerate}
  \item The imaginary sector $h_{\mu\nu}$ is non-trivial and provides
    negative effective energy density in a region $\mathcal{V}$.
  \item The real metric $g_{\mu\nu}$ in $\mathcal{V}$ contains a Killing
    vector $\xi^\mu$ with $g_{\mu\nu}\xi^\mu\xi^\nu < 0$ and compact orbits.
  \item Either (a) $\psi$ is a second independent time direction
    (Interpretation A), or (b) the energy condition violation from
    $h_{\mu\nu}$ is large enough to tilt light cones in the real sector.
  \item The Cauchy-horizon back-reaction from the UBT quantum stress-energy
    is finite (the KK sum \eqref{eq:kk_sum} in ST-4 converges).
\end{enumerate}
\end{conjecture}

\begin{remark}
The conjunction of conditions~1--4 is not yet known to be sufficient.
Additional conditions may be required, for example topological
non-triviality of the $\psi$-fibre bundle or specific boundary conditions
on the $\Theta$-field at spatial infinity.
\end{remark}

%----------------------------------------------------------------------
\section{Minimal Modification to Canonical Field Equations}
%----------------------------------------------------------------------
The canonical UBT field equation is
\begin{equation}
  \mathcal{E}_{\mu\nu}[\Theta,\mathcal{G}] \;=\;
  \kappa\,\mathcal{T}_{\mu\nu}[\Theta,\mathcal{G}].
  \label{eq:canonical_ubt}
\end{equation}
This equation, in the real projection, recovers Einstein's equations with
an effective stress-energy sourced by $\Theta$.  For CTCs to be
\emph{dynamically stable} solutions (rather than fine-tuned initial data),
the field equation must actively support negative energy density regions.

The minimal modification we identify is the addition of a
\emph{complex cosmological term}:
\begin{equation}
  \mathcal{E}_{\mu\nu}[\Theta,\mathcal{G}]
  + \Lambda_\mathbb{B}\,\mathcal{G}_{\mu\nu}
  \;=\;
  \kappa\,\mathcal{T}_{\mu\nu}[\Theta,\mathcal{G}],
  \label{eq:modified_ubt}
\end{equation}
where $\Lambda_\mathbb{B} = \Lambda_r + i\Lambda_i$ is a
\emph{biquaternionic cosmological constant}.  The imaginary part
$\Lambda_i$ contributes directly to the effective energy density in the
imaginary sector:
\begin{equation}
  \rho_h^{(\Lambda)} \;=\; -\frac{\Lambda_i}{8\pi G}.
\end{equation}
For $\Lambda_i > 0$ this generates positive imaginary energy density,
which in the effective real equation \eqref{eq:einstein_eff} acts as a
negative energy density component — precisely the energy condition
violation required for CTCs.

\begin{conjecture}[Minimal CTC-enabling modification]
\label{conj:minimal_mod}
\textbf{(Speculative)}  The addition of a biquaternionic cosmological
constant $\Lambda_\mathbb{B}$ with $\operatorname{Im}(\Lambda_\mathbb{B})
> 0$ to the canonical field equation \eqref{eq:canonical_ubt} constitutes
the minimal modification that enables CTC-compatible solutions.  No
modification of the kinetic or interaction terms of $\Theta$ is required.
\end{conjecture}

\begin{remark}
This conjecture does not assert that such a modification is physically
motivated or that $\Lambda_\mathbb{B} \neq 0$ is required by the canonical
UBT framework.  It identifies the minimal mathematical change.
\end{remark}

%----------------------------------------------------------------------
\section{Time-Travel Protocol in UBT Field Language}
%----------------------------------------------------------------------
A ``time-travel protocol'' in UBT refers to a sequence of field
configurations of $\Theta$ that, if realised, would correspond to a
CTC in the spacetime geometry.  We describe such a protocol
\emph{in principle} — not as engineering instructions, but as a
characterisation of the field configurations involved.

\paragraph{Step 1: Establish a rotating phase configuration.}
Configure $\Theta$ in a spatially localised, rotating state:
\begin{equation}
  \Theta_0(q,\tau) \;=\; A_0\,e^{i(\omega_\psi\psi - \omega_t t)}\,
  f(\mathbf{x}),
\end{equation}
where $f(\mathbf{x})$ is a localised envelope function and
$\omega_\psi/\omega_t = r_\omega$ is the ratio of phase rotation to real
rotation rates.

\paragraph{Step 2: Seed imaginary metric sector.}
The rotating $\Theta_0$ sources a non-zero $h_{\mu\nu}$ through the
imaginary part of the biquaternionic stress-energy:
\begin{equation}
  h_{\mu\nu}^{(\mathrm{source})} \;\propto\;
  \operatorname{Im}\!\left\langle D_\mu\Theta_0, D_\nu\Theta_0
  \right\rangle_\mathbb{B}.
\end{equation}
For $r_\omega > r_\omega^{\mathrm{crit}}$ (a critical ratio determined by
the UBT coupling constants), $h_{\mu\nu}$ attains a magnitude sufficient
to locally violate the WEC via the effective stress-energy.

\paragraph{Step 3: Develop CTC-compatible real metric.}
Back-reaction of $h_{\mu\nu}$ on $g_{\mu\nu}$ through equation
\eqref{eq:modified_ubt} deforms the causal structure of the real sector.
If conditions 1--4 of Conjecture~\ref{conj:necessary} are satisfied, the
real metric develops a region with a timelike Killing vector with compact
orbits — a CTC region.

\paragraph{Step 4: Persistence or dissolution.}
The CTC region either:
\begin{itemize}
  \item \textbf{Persists}: if the KK sum converges (Scenario B of ST-4),
    the Cauchy-horizon back-reaction is finite and the CTC is stable.
  \item \textbf{Dissolves}: if the KK sum diverges (Scenario A of ST-4),
    the back-reaction destroys the CTC region and chronology protection
    is enforced.
\end{itemize}

\begin{tcolorbox}[colback=orange!10, title=Critical Framing Note]
The above protocol describes what field configurations $\Theta$ would be
\emph{involved} if CTCs exist in UBT.  It does not claim that:
\begin{itemize}
  \item Such configurations can be realised by any physical process.
  \item The energy scales involved are accessible to any known technology.
  \item The UBT framework permits CTC formation.
  \item Time travel is physically possible.
\end{itemize}
The protocol is a \emph{theoretical characterisation}, not a claim about
physical accessibility.
\end{tcolorbox}

%----------------------------------------------------------------------
\section{Summary: Open Questions and Research Directions}
%----------------------------------------------------------------------
The analysis of ST-1 through ST-4 identifies the following open questions
as the critical bottlenecks for establishing whether UBT admits CTCs:

\begin{enumerate}
  \item \textbf{Interpretation of $\psi$:}  Resolve whether $\psi$ is an
    internal phase (Interpretation B) or a second time direction
    (Interpretation A) by deriving the propagator of the $\Theta$-field
    in both cases and comparing with observable predictions.
  \item \textbf{Convergence of the KK sum:}  Evaluate the Kaluza–Klein
    sum of the renormalised stress-energy on a prototypical Cauchy horizon
    in UBT.  This requires the full UBT propagator on a CTC background.
  \item \textbf{Stability of imaginary energy solutions:}  Identify
    explicit solutions of \eqref{eq:canonical_ubt} with $h_{\mu\nu} \neq 0$
    and check whether they satisfy the CTC conditions of ST-3.
  \item \textbf{$U(1)_\tau$ symmetry:}  Verify or refute the existence
    of the temporal $U(1)$ symmetry (Conjecture~2 of ST-2) by computing
    the variation of the UBT action under $\tau \mapsto e^{i\alpha}\tau$.
\end{enumerate}

\paragraph{Null result.}
If the above analysis ultimately shows that UBT does not admit CTCs more
naturally than GR — for example because the KK sum always diverges or
because $\psi$ is strictly an internal phase — that is an equally valid
and publishable result, establishing a rigorous ``no-CTC theorem'' for
UBT.

\medskip
\noindent\textbf{All statements in this document are speculative.  No
claim constitutes an established result of the canonical UBT framework.
No conjecture may be promoted to a Theorem or Lemma without independent
proof in the canonical framework.}

%----------------------------------------------------------------------
\bibliographystyle{plain}
\begin{thebibliography}{9}
\bibitem{Hawking1992}
  S.~W.~Hawking,
  \emph{Chronology protection conjecture},
  Phys.\ Rev.\ D \textbf{46}, 603 (1992).
\bibitem{Godel1949}
  K.~Gödel,
  \emph{An example of a new type of cosmological solution of Einstein's
  field equations of gravitation},
  Rev.\ Mod.\ Phys.\ \textbf{21}, 447 (1949).
\bibitem{UBT_core}
  D.~Jaroš,
  \emph{Unified Biquaternion Theory: core field equations},
  Repository document \texttt{consolidation\_project/ubt\_core\_main.tex} (2025).
\bibitem{CTC_ST1}
  D.~Jaroš,
  \emph{ST-1: Causal Structure of Complex Time},
  \texttt{speculative\_extensions/causality/complex\_time\_causal\_structure.tex}
  (2026).
\bibitem{CTC_ST2}
  D.~Jaroš,
  \emph{ST-2: Rotation in the $(t,\psi)$ Plane},
  \texttt{speculative\_extensions/causality/rotation\_complex\_time.tex}
  (2026).
\bibitem{CTC_ST3}
  D.~Jaroš,
  \emph{ST-3: CTC Conditions in the Biquaternionic Metric},
  \texttt{speculative\_extensions/causality/ctc\_conditions.tex} (2026).
\bibitem{CTC_ST4}
  D.~Jaroš,
  \emph{ST-4: Chronology Protection in UBT},
  \texttt{speculative\_extensions/causality/chronology\_protection\_ubt.tex}
  (2026).
\end{thebibliography}

\end{document}
